\section{Related Work}
\label{sec:related}

\paragraph{Aggregated deletion propagation.}
Hu et al.~\cite{hu2021adp} study minimum source deletions to remove a prescribed
number of query outputs. In the shared
setting of full self-join-free queries over duplicate-free relations,
each full output corresponds to one derivation, so their ADP objective
is exactly our \agg{COUNT} objective. Their Theorem~6.2 gives
$O(\log k)$ and relation-count approximations, using partial set cover.
Bag projection here counts derivation copies;
set projection coalesces them and can change marginal behavior, as the
example in Section~\ref{sec:problem} shows. Weighted values, designated deletability, and aggregate validity conditions
require comparison under their actual semantics.

% Balance the final page with the class-provided balance package.
\balance
\paragraph{Generalized deletion propagation and resilience.}
Makhija and Gatterbauer~\cite{makhija2025ilp} unify deletion and preservation
requirements with source and view objectives. Their framework handles
self-joins and bags and gives ILP formulations with strong relaxation
properties in known tractable cases. Our survivor constraints and tuple-to-
derivation linking belong in this established methodological context.
The aggregate-specific inequalities in Section~\ref{sec:ilp} make our
threshold semantics explicit; we do not assert that this encoding is
outside their framework or has the same relaxation properties.

At $\theta=Q(D)>0$, \agg{COUNT} requires destroying all derivations,
which is query resilience with the same permitted deletions. Freire
et al.~\cite{freire2015resilience} study resilience as minimum deletions
making a Boolean conjunctive query false and characterize self-join-free
query classes. Our reduction uses a self-join and immutability; it neither
replaces that classification nor extends it to all aggregate queries.

\paragraph{Provenance and aggregate semantics.}
Semiring provenance represents the derivational dependence of positive
queries~\cite{green2007provenance}; aggregate provenance develops the
corresponding treatment of aggregate values~\cite{amsterdamer2011aggregate}.
We use fixed derivation values and deletable supports to reason about
interventions on that dependence. Our Boolean circuit contract is a
specific operational requirement for exact marginal computation, rather
than a claim that every provenance representation supplies the same oracle.
The multilinearity qualification is essential when tuple identifiers recur.

\paragraph{Submodular cover and bounded-rank hitting set.}
The logarithmic argument specializes the standard greedy analysis for
submodular covering~\cite{wolsey1982greedy}. The disjoint-support argument
is the elementary packing bound for bounded-rank hitting set. Our extrema
analysis uses these ingredients after selecting a survivor, so that an
otherwise feasible critical cover does not empty the aggregate input.
This interaction between threshold obligations and preservation matters
when applying deletion-propagation algorithms to aggregates that exclude
empty inputs.
